\documentclass{article}
\usepackage{algorithm, algorithmic, amsfonts, amsmath}

\begin{document}
\begin{algorithm}[!h]
\renewcommand{\algorithmicrequire}{\textbf{Input:}} % show \REQUIRE as Input
\renewcommand{\algorithmicensure}{\textbf{Output:}} % show \ENSURE as Output
	\caption{PULNS (reinforced negative sample selection)} % algorithm title
	\begin{algorithmic}[1] % [1] numbers every line
	    \REQUIRE
	    Positive set $P$; unlabeled set $U$; validation set $V$ (contains a \emph{small} number of
	    labeled negatives); episode number $L$; reward weight $\alpha$
	    \ENSURE
	    Optimized selector $f_{\theta}^{\ast}$; optimized classifier $clf^{\ast}$
	    \STATE Initialize a classifier $clf$ trained on $P$ and $U$, treating all of $U$ as negative % a biased starting point, corrected by the iterations that follow
	    \STATE Initialize a selector $f_{\theta}$ randomly
	    \STATE $f_{\theta}^{\ast}\leftarrow f_{\theta}$;\quad $clf^{\ast}\leftarrow clf$
	    \FOR{$l\leftarrow 1$ \TO $L$}
	    \STATE Shuffle the unlabeled sample set $U$
	    \FOR{step $i\leftarrow 1$ \TO $|U|$}
	    \STATE Get state $s_{i}$
	    \STATE Sample action $a_{i}\sim \pi_{\theta}(s_{i},a_{i})$ \quad ($a_{i}=1$: select as negative) % the selector decides instance by instance whether to take it
	    \STATE Obtain intermediate reward
	    $r_{i}^{(l)}=\begin{cases}
	    C\!\left(-1,\ \log\frac{1-p(x_{i})}{p(x_{i})},\ 1\right) & \text{if } a_{i}=1\\[2pt]
	    C\!\left(-1,\ \log\frac{p(x_{i})}{1-p(x_{i})},\ 1\right) & \text{if } a_{i}=0
	    \end{cases}$ % C(-1,x,1) is a clipping function; the more negative the classifier leans, the higher the reward for selecting it
	    \ENDFOR
	    \STATE Compute terminal reward $r_{u+1}$ from the classifier's performance on $V$ % the terminal reward comes from a validation set, so a few true negatives are needed
	    \STATE Update selector $f_{\theta}$ based on all rewards:
	    $\nabla_{\theta}J(\theta)=\sum_{\tau}\nabla_{\theta}P_{\theta}(\tau)\left(\sum_{i=1}^{u}r_{i}+\alpha r_{u+1}\right)$ %REINFORCE
	    \STATE $f_{\theta}^{\ast}\leftarrow f_{\theta}$
	    \STATE Get calibrated negative sample set $N$ from $U$ by the updated selector $f_{\theta}^{\ast}$
	    \STATE Update classifier $clf$ based on $N$ and $P$ % selector and classifier are optimized alternately
	    \IF{$clf$ performs better than $clf^{\ast}$}
	    \STATE $clf^{\ast}\leftarrow clf$
	    \ENDIF
	    \ENDFOR
		\RETURN $f_{\theta}^{\ast}$, $clf^{\ast}$
	\end{algorithmic}
\end{algorithm}
% Luo et al., AAAI 2021, "PULNS: Positive-Unlabeled Learning with Effective Negative
% Sample Selector". Compiled from Algorithm 1 and the intermediate reward definitions of the original.
% Difference from the classical two-step methods: those select negatives once by a heuristic rule (spies, Rocchio,
% 1-DNF and so on), whereas PULNS models the choice of negatives as a sequential decision problem, learns a selector
% by REINFORCE with classifier performance as the reward, and optimizes the two alternately.
% The reward has two levels:
% (1) an intermediate reward per instance, which avoids the sparse return of a terminal reward alone;
% (2) a terminal reward from validation performance. The original notes that using accuracy directly impedes
%     convergence, so a different form is used. Prerequisite: a validation set holding a few true negatives. The
% original argues this is practical, since collecting negatives at scale is unrealistic but labelling a few for
% validation usually is, and per Hsieh et al. (2019) even a biased sample of them still yields a sizable gain.
\end{document}
